\section{Two stages recover the original sparse-source RPPR objective}
\label{sec:revisit-general-rppr}
The semantic source truncation has a stronger two-stage consequence.
Its first stage changes both source and regularizer to find an approximate
coordinate envelope. The second stage restores the original source and the
requested regularizer. Thus it solves the original RPPR objective rather
than treating the source change as part of the final error definition.

\begin{lemma}[Source reduction compensated by lower regularization]
\label{lem:revisit-source-envelope}
Fix $\rho>0$. Retain $T_{\rm h}=\{i:s_i/d_i>\rho/2\}$, let
$\vs_{\rm h}$ contain their original weights, and let $\vu_{\rm h}$
be its RPPR optimum at regularizer $\rho/2$. Then
\begin{equation}
 \vx_\rho^*(\vs)\leq\vu_{\rm h},\qquad
 \vol(T_{\rm h})<2/\rho,\qquad
 \vol(\supp\vu_{\rm h})\leq2/\rho.
 \label{eq:revisit-source-envelope}
\end{equation}
\end{lemma}
\begin{proof}
The omitted load satisfies
$0\leq\vb_{\rm l}\leq(\alpha\rho/2)\vomega$.
Consequently
\[
 \vb_{\rm h}-(\alpha\rho/2)\vomega
 \geq\vb-\alpha\rho\vomega.
\]
Solutions of an orthant quadratic with a Stieltjes Hessian are monotone
in this linear load. To verify the comparison directly, let $\va,\vb'$ be
the optima for loads $\vc_a\geq\vc_b$, and put
$\vd=(\vb'-\va)_+$. On its positive support, complementarity for
$\vb'$ and the KKT inequality for $\va$ give
$[\mQ(\vb'-\va)]_i\leq0$.
Replacing the negative coordinates of $\vb'-\va$ by zero can only decrease
this product on that support. Hence $\vd^\trans\mQ\vd\leq0$,
which forces $\vd=0$.
Apply this fact to the displayed loads. The two volume bounds follow
respectively by summing retained source weights and by RPPR residual mass.
\end{proof}

\begin{theorem}[Sparse-source RPPR with additive input work]
\label{thm:revisit-general-rppr}
Given an explicit sparse source list of size $|T|$, $\alpha\in(0,1]$,
$\rho>0$, and $0<\epsobj<1$, there is a deterministic two-stage local
algorithm returning
\[
 0\leq\widehat\vx\leq\vx_\rho^*(\vs),\qquad
 \vol(\supp\widehat\vx)\leq1/\rho,
 \qquad F_\rho(\widehat\vx)-F_\rho(\vx_\rho^*)\leq\epsobj
\]
in fully charged exact-real work
\[
 \widetilde O\!\left(|T|+\frac1{\rho\sqrt\alpha}\right),
\]
with logarithmic accuracy dependence. A separately rounded realization
retains this local operation scale and bounded integer encodings.
\end{theorem}
\begin{proof}[Exact-real construction and accuracy]
Read the source list and its degrees. If
$\rho\geq\max_i s_i/d_i$, return zero; the KKT conditions hold there.
At $\alpha=1$ return the exact masses $[s_i-\rho d_i]_+$.
In the remaining case $\rho<1$, retain $T_{\rm h}$ as above, choose a
dyadic $\delta$ by halving $1/2$ until
$\delta^2\leq\epsobj\rho/4$, and proceed as follows.

Stage I runs the joint retained-source orthant continuation at final
regularizer $\rho/2$. Freeze it until its actual Euclidean accelerated
energy is at most $\alpha\delta^2/2$, and set
\[
 \vp=[\vx-\delta\vomega]_+,\qquad U=\supp\vp.
\]
The subprobability-source proof in \cref{sec:revisit-sparse-sources} gives
$0\leq\vp\leq\vu_{\rm h}$ and
$0\leq\vu_{\rm h}-\vp\leq2\delta\vomega$.
Thus $\vol(U)\leq2/\rho$, and
\cref{lem:revisit-source-envelope} gives
\begin{equation}
 i\notin U\quad\Longrightarrow\quad
 (\vx_\rho^*)_i\leq2\delta\omega_i.
 \label{eq:revisit-missing-source-support}
\end{equation}
If $U$ is empty, the same argument below certifies returning zero.

Stage II builds the principal matrix $\mQ_{UU}$ using original degrees
and the \emph{original} source coefficients on $U$. Let $\vu_U$ be the
exact restricted RPPR optimum, zero-extended. It is below $\vx_\rho^*$:
the latter's restriction to $U$ is a supersolution for the principal
obstacle problem, and the Stieltjes comparison proof above applies.
Run fixed-face orthant acceleration from zero until its energy is at most
$\alpha\delta^2/2$, then clip by $\delta\vomega_U$ and zero-extend.
The initial energy is at most one because
$\|\vu_U\|_2\leq\vomega^\trans\vu_U\leq1$.
This returns $0\leq\widehat\vx\leq\vu_U$ and degree-density error at
most $2\delta$ from $\vu_U$.

To bound the envelope error, truncate the true full optimum to $U$.
The removed vector is supported on $\supp\vx_\rho^*$ and is bounded
by $2\delta\vomega$, by \cref{eq:revisit-missing-source-support}.
The KKT linear term vanishes on that optimal support, so
\[
 F_\rho(\vu_U)-F_\rho(\vx_\rho^*)
 \leq F_\rho((\vx_\rho^*)_U)-F_\rho(\vx_\rho^*)
 \leq2\delta^2/\rho.
\]
The same quadratic expansion for the final clipped error, supported on
$\supp\vu_U\subseteq\supp\vx_\rho^*$, gives
$F_\rho(\widehat\vx)-F_\rho(\vu_U)\leq2\delta^2/\rho$.
Their sum is at most $\epsobj$. Lower order also proves the final support
volume bound.
\end{proof}

\begin{proof}[Work and bounded arithmetic]
The initial source pass costs $O(|T|)$ degree replies and input reads.
Stage I refreshes only $T_{\rm h}$ as source exceptions, of cardinality at
most $2/\rho$. Its iteration count is
$O(\alpha^{-1/2}\log(2/(\rho\epsobj)))$, and its kinetic work has the
same $44\sum_k1/r_k$ bound (or $60\sum_k1/r_k$ when rounded).
Source-exception updates therefore fit the claimed total.
No source row is opened merely because its label was supplied in the input.

Building the principal face costs $O(\vol(U))$ incidence work plus indexing
logarithms. Retrieve its original source entries by one further pass over
the supplied source list and membership tests in $U$; this costs
$\widetilde O(|T|)$ and does not query cold-source adjacency lists.
Stage II makes
$O(\alpha^{-1/2}\log(2/(\alpha\rho\epsobj)))$ fixed-face iterations,
each costing $O(\vol(U))$ with all repeated original incidences charged.
Final clipping and output add $O(1+\vol(U))$ scalar work.

The bounded Stage I uses \cref{sec:revisit-rounding}, with its frozen
tolerance replaced by $\alpha\delta^2/2$. For Stage II let
$V_U=\vol(U)$, use a dyadic density grid
\[
 h_U\leq\frac{\delta\theta^2}{16\max\{1,V_U\}},
\]
and round both kinetic and next primal densities downward after each dense
fixed-face step. Choose $N=q/\theta$, where $q$ is the least integer with
$2^q\geq8/(\alpha\delta^2)$.
The corresponding exact trajectory has Euclidean error at most $\delta/2$.
Each rounded step differs from its ideal step by at most $h_U$ in kinetic
density and $(1+\theta)h_U$ in primal density.
The pairwise-contraction proof of \cref{lem:revisit-pairwise} applies to
the principal matrix, giving accumulated state-distance error at most
$4h_U\sqrt{V_U}/\theta$ and hence primal Euclidean error at most
\[
 \frac{4\sqrt2h_U\sqrt{V_U}}{\theta\sqrt\alpha}<\delta/2.
\]
Thus the actual candidate is within $\delta$ of $\vu_U$, exactly what the
clipping proof requires.

Every new Stage II state returns immediately to its fixed grid. Exact
source terms use only individual input denominators, and all neighbor sums
are sums of grid integers. The exact trajectory's initial energy is at
most one, which bounds its primal and kinetic norms polynomially in
$1/\alpha$; the comparison bound gives the same for the rounded state.
Integer lengths are therefore bounded by input encoding lengths and
logarithms of $1/(\alpha\rho\epsobj)$ and encountered degrees.
No denominator grows linearly with the iteration count. This finishes
the separately rounded construction.
\end{proof}

\begin{lemma}[Maximum principle for an incomplete obstacle face]
\label{lem:revisit-obstacle-face-error}
Let $\vu$ be the full RPPR optimum and $\vu_U$ its restricted optimum
on $U$, zero-extended. If $u_i\leq a\omega_i$ outside $U$, then
\[
 \vzero\leq\vu-\vu_U\leq a\vomega.
\]
\end{lemma}
\begin{proof}
The lower bound is the Stieltjes comparison already used above. Write
$\vc=\vb-\alpha\rho\vomega$ and $\mM=\mQ_{UU}$.
The full optimum restricted to $U$ solves the obstacle problem with load
$\vc_U-\mQ_{U,U^c}\vu_{U^c}$. Meanwhile
$\vu_{U,U}+a\vomega_U$ is a supersolution, because
\[
 \mM(\vu_{U,U}+a\vomega_U)-\vc_U
                +\mQ_{U,U^c}\vu_{U^c}
 \geq a(\mM\vomega_U+\mQ_{U,U^c}\vomega_{U^c})
 =a\alpha\vomega_U\geq\vzero.
\]
Here nonpositive off-diagonals and the exterior amplitude bound give the
inequality. The same positive-part comparison with a supersolution proves
the claimed upper bound on $U$; it is assumed on the complement.
\end{proof}

\begin{corollary}[General-source residual and semantic certificates]
\label{cor:revisit-general-residual}
Run the two-stage construction with any dyadic $\delta>0$, using its
stated $\alpha\delta^2/2$ accuracy budget in Stage I and the same Stage II
schedule. Then
\[
 \vzero\leq\vx_\rho^*-\widehat\vx\leq4\delta\vomega,
 \qquad
 \|\mD^{-1/2}\bm{\kappa}(\widehat\vx)\|_\infty\leq4\delta,
\]
where $\bm{\kappa}$ is the minimum-magnitude RPPR subgradient.
The semantic PPR error is at most $\rho+4\delta$.
If $\delta\leq\alpha\rho/4$, the full original-source residual additionally
satisfies
\[
 \vzero\leq\vb-\mQ\widehat\vx\leq2\alpha\rho\vomega.
\]
All work bounds retain their scale, with logarithmic dependence on
$1/\delta$.
\end{corollary}
\begin{proof}
Stage I bounds omitted optimal coordinates by $2\delta\vomega$.
Apply \cref{lem:revisit-obstacle-face-error}, then add the Stage II clipping
error of at most $2\delta\vomega$. Write their total as
$\ve=\vx_\rho^*-\widehat\vx$, so
$|\mQ\ve|\leq4\delta\vomega$, using $|\mQ|\vomega\leq\vomega$.
At every positive output coordinate the exact optimum has zero RPPR
gradient. At zero output coordinates only residual excess over
$\alpha\rho\vomega$ contributes to the minimum subgradient. These facts
give the subgradient bound, and the RPPR bias gives the semantic bound.
For residual nonnegativity, positive output coordinates have residual at
least $(\alpha\rho-4\delta)\vomega$. At a zero output coordinate,
Stieltjes signs give nonnegative residual directly. The residual upper
bound follows from $\vb-\mQ\vx_\rho^*\leq\alpha\rho\vomega$.
\end{proof}

The second stage is needed to restore the original RPPR objective. In
general, the first-stage optimum has a different source and a smaller
regularizer, so directly returning it would not satisfy the requested
objective-gap statement. Neither stage needs exact support identification.
